\documentclass{../template/note}

\author{Oliver West}
\title{Dynamics and relativity}
\date{}

\begin{document}

    \maketitle

    \tableofcontents

    \lecture{Lecture 1 -- 22/01/26}

    \section{The structure of the Newtonian universe}

    \begin{itemize}
        \item Three dimensional space endowed with a \emph{Cartesian reference frame} -- an origin and some axes.
        \item Time can be labelled in an arbitrary reference frame, by a real number $t$.
        \item A \emph{point particle} is an idealised object that is completely determined by its position at a given point in time, $\vec x(t)$. Its velocity
        \begin{align*}
            \vec v = \frac{d \vec x}{dt} = \left(\frac{dx_1}{dt}, \frac{dx_2}{dt}, \frac{dx_3}{dt}\right) = \dot{\vec x}
        \end{align*}
        is tangent to the trajectory. The acceleration
        \begin{align*}
            \vec a = \frac{d^2\vec x}{dt^2} = \left(\frac{d^2x_1}{dt^2}, \frac{d^2x_2}{dt^2}, \frac{d^2x_3}{dt^2}\right) = \ddot{\vec x} = \dot{\vec v}
        \end{align*}
    \end{itemize}
    The above is \emph{not} sufficient to write down Newton's equations. Consider, for example, a free particle not experiencing any forces. The position of this particle is $\vec x(t)$, but which frame should we use?
    
    \subsection{Law of inertia}

    \begin{proposition}[Law of inertia]
        There exist \emph{inertial frames} in which a free particle has a constant velocity.
    \end{proposition}
    
    This is an improved version of Newton's first law. Indeed, it is a true statement, but not an obvious one.

    \subsection{Galilean relativity principle}
    
    \begin{definition}
        A \emph{Galilean transformation} has the form
        \begin{align*}
            \vec x' = R \vec x + \vec k + \vec w t
        \end{align*}
        where $R \in SO(3)$, $\vec k$ is a translation vector, and $\vec w$ is a \emph{boost}. In particular, $\ddot{\vec x'} = \ddot{\vec x}$.
    \end{definition}

    \begin{proposition}[Galilean relativity principle]
        A frame related to an inertial frame by a \emph{Galilean transformation} is also an inertial frame, and all laws of physics are the same in both frames.
    \end{proposition}
    
    Example of a boost: dropping a ball on a boat moving at constant velocity.

    Galilean invariance also restricts the forces that are possible (see example sheet). For example, there are no `special points' in the universe (or special directions, times, velocities, etc) -- they are all relative. This is \emph{not} the case for acceleration -- an object accelerating in one inertial frame is accelerating with the same magnitude in all inertial frames.

    Galilean transformations form the \emph{Galilean group}, often supplemented by shifts in time, under which the laws of physics are also invariant. In particular, this implies that all inertial frames have the same time, called absolute time.

    \section{Forces}

    Once there is more than one particle in the universe, there will be interactions between them. In Newtonian physics, these are described by forces.

    \subsection{Newton's second law}

    In an inertial frame,

    \begin{align*}
        \dot{\vec p} = \vec F
    \end{align*}
    
    where the \emph{momentum} $\vec p$ is $m \vec v$, where $m$ is the \emph{inertial mass}, mass for short. The mass is an additional property for point particles that we take to be constant in time unless otherwise stated.

    The force $\vec F$ depends on the ongoing interactions, but it can only depend on $\vec x$ and $\vec v$. This implies that Newton's second law is a second-order differential equation, and that, given $\vec x$ and $\vec v$ at $t=0$ for all particles, Newton's equations will uniquely determine $\vec x(t)$ for all future times.

    We should note that Newtonian mechanics has been superseded by quantum mechanics (for small things) and general relativity (for fast things).
    
    
    
    
    
    
\end{document}
